\section{Ausführung mehrerer Dateien} \label{sec:multi_file_execution}
In diesem Abschnitt wird die Ausführung des Algorithmus auf mehreren GeoTIFF-Dateien beschrieben. Dabei wird erläutert, wie die bestehenede Architektur verändert werden muss um eine Verarbeitung mehrerer Dateien zu ermöglichen. Dabei sit zu beachten, dass die Dateien zunächst sequenziell verabreitet werden sollen. Die parallele Verarbeitung von Dateien wird in \ref{sec:multi_file_parallel} beschrieben. Außerdem soll die Möglichkeit der Verabreitung einer einzelnen Datei mit dem alten Algorithmus weiterhin bestehen bleiben. Dies ist dem geschuldet, dass manche Kartenabschnitte nur in einer Datei vorliegen. 

Als ersten Schritt muss die Möglichkeit geschaffen werden, mehrere Dateien als Eingabe zu akzeptieren. Dafür muss in dem Tester aus \ref{chap:bewertung_gipfelfindungsalgorithmus} eine neue Variable namens \textit{input\_folder} erstellt werden. Diese Variable definiert den Pfad zu einem Ordner, in dem sich die GeoTIFF-Dateien befinden, die verarbeitet werden sollen. Diese Varibale darf nur gesetzt sein, wenn \textit{input\_file} nicht gesetzt ist. Außerdem muss immer eine der beiden varibalen gestzt sein. Es muss eine Logik implementiert werden, die alle GeoTIFF-Dateien in diesem Ordner identifiziert und in einer Liste speichert. Diese Liste wird dann verwendet, um die Dateien nacheinander zu verarbeiten.

Die Verarbeitung der Dateien erfolgt in einer Schleife, die über die Liste der Dateien iteriert. Für jede Datei wird der Algorithmus wie gewohnt ausgeführt, um die Gipfel zu identifizieren. Diese Funktion gibt für jede Datei eine Liste von Gipfeln zurück, die in dieser Datei gefunden wurden.  Diese Gipfel sind jedoch nur lokal für die jeweilige Datei gültig. Um zu bestimmen, in wie fern die bestimmten Gipfel über die Dateigrenzen hinweg gültig sind, wird in \ref{sec:multi_file_merge} beschrieben, wie die Prominenz- und Dominanzberechnung über die Dateigrenzen hinweg ermöglicht wird.

Bezüglich des Randwertproblems muss beachtet werden, dass die Gipfel, die sich im Randwert einer Datei befinden nach der aktuellen Implementierung ignoriert werden, obwohl ausreichend umgebende Daten vorhanden sind um die Prominenz- und Dominanzberechnung später dateiübergreifend auszuführen. Deshalb wird das Randwertproblem in diesem Fall nicht auf der Einzelfallebene gelöst, sondern erst bei der Zusammenführung der Ergebnisse über die Dateigrenzen hinweg. Das bedeutet, dass die Gipfel, die sich im Randbereich einer Datei befinden, zunächst als gültige Gipfel betrachtet werden. Erst bei der Zusammenführung der Ergebnisse über die Dateigrenzen hinweg wird überprüft, ob diese Gipfel tatsächlich gültig sind oder ob sie aufgrund von Daten in benachbarten Dateien als ungültig betrachtet werden müssen. Diese Vorgehensweise ermöglicht es, die Gipfel, die sich im Randwert einer Datei befinden, nicht vorschnell zu verwerfen und somit potenziell gültige Gipfel zu erhalten, die sonst verloren gehen würden. 

\todo{file analysis manager class}